% Generated draft: requires mathematical and source review.
\documentclass[11pt]{article}
\usepackage[margin=25mm]{geometry}
\usepackage{fontspec}
\setmainfont{Latin Modern Roman}
\usepackage{amsmath,amssymb,mathtools,unicode-math}
\setmathfont{Latin Modern Math}
\usepackage{enumitem,needspace,xurl,hyperref}
\setlength{\parindent}{0pt}
\setlength{\parskip}{5pt}
\setlength{\emergencystretch}{3em}
\newcommand{\N}{\mathbb{N}}
\newcommand{\Z}{\mathbb{Z}}
\newcommand{\Q}{\mathbb{Q}}
\newcommand{\R}{\mathbb{R}}
\newcommand{\C}{\mathbb{C}}
\newcommand{\F}{\mathbb{F}}
\newcommand{\A}{\mathbb{A}}
\newcommand{\PP}{\mathbb{P}}
\newcommand{\E}{\mathbb{E}}
\newcommand{\eps}{\varepsilon}
\newcommand{\dd}{\,\mathrm{d}}
\begin{document}
\section*{UnsolvedMath ID: 30000873}
\textbf{Source problem code:} catalogue entry
\subsection{Definitions and mathematical statement}
This record poses a mathematics-education research agenda about mathematical knowledge for teaching, rather than a mathematical conjecture with specified objects and quantifiers. The expression concerns the mathematical knowledge teachers possess or require and their ability to use it in teaching; the source does not select a particular theoretical framework or operational measure. Preserve the following linked questions:
\begin{enumerate}
\item What perspectives on mathematical knowledge for teaching have been proposed, and what theoretical or empirical bases support each perspective?
\item What methods have been used to study what teachers know, what they use, what they need to know, and what they need to know how to use? These four aspects are distinct in the source.
\item Which aspects of those methods are common and which are distinct, and how have the methods shaped the research questions being asked?
\item How can quantitative, qualitative, and interpretive results support teachers' mathematical knowledge and their professional views of mathematics in teaching?
\end{enumerate}
The requested work is an account and comparison of frameworks, evidential bases, research methods, and implications for professional development. It is not specified as the existence of one final framework, a universal causal effect, or a numerical optimization problem. A formal empirical study would additionally need a teacher population, educational and cultural settings, constructs, instruments, outcomes, and criteria for evaluating support; these are not supplied. No artificial mathematical theorem is substituted for the agenda.
\subsection{Short English statement}
Compare perspectives and evidence about mathematical knowledge for teaching, the methods used to study teachers' knowledge and its use, how those methods shape research, and how the results can support teachers' knowledge and professional views.
\subsection{Sources}
\begin{itemize}
\item Local UnsolvedMath registry, permanent ID \(30000873\), source code catalogue entry, ``Frameworks for Mathematical Knowledge in Teaching.'' No individual proposer is supplied.
\item \emph{Professional Development of Mathematics Teachers---Research and Practice from an International Perspective} (2008), Oberwolfach Reports, DOI 10.4171/owr/2007/52, pp. 3073--3130: \url{https://doi.org/10.4171/owr/2007/52}.
\item Supplied reference, ``international perspective review'': \url{https://doi.org/10.1016/j.jmathb.2018.05.002}.
\item Supplied reference, ``comparative MKT study'': \url{https://link.springer.com/article/10.1007/s10857-023-09598-z}.
\item Supplied literature triage, checked 2026-08-21: status ``open.'' It reports multiple competing theoretical and empirical frameworks and no established final integrating account. This is a dated assessment derived from supplied OWR curation, not an independently verified empirical or mathematical conclusion; time-sensitive status and cited claims require independent verification.
\end{itemize}
\end{document}
